\section{Methods}\label{sec:methods}

Figure~\ref{fig:pipeline} summarises the protocol. Players were split once, at random within tertiles of career length, into a development set (80\%, $n=\nDev$) and a held-out set (20\%, $n=\nHeld$). The development set was used for autoencoder training, standardisation constants and every clustering hyper-parameter choice; the held-out set for early stopping of the autoencoder, for reporting reconstruction error and for testing the transferability of clustering solutions. Because the same 20\% served early stopping and reporting, the reported reconstruction error is slightly optimistic; a 70/10/20 check (Appendix~\ref{app:training}) bounds the effect. The consensus partitions interpreted in Section~\ref{sec:results} were computed on all players after these checks.

\begin{figure}[t]\centering
\includegraphics[width=\textwidth]{fig2_pipeline}
\caption{Analytical protocol. Top: data construction and the player-level split. Middle: representations compared on equal terms, and the single-run design of the earlier version, evaluated for reproducibility. Bottom: the common consensus protocol and the tests for discrete structure.}
\label{fig:pipeline}
\end{figure}

\subsection{Transformer autoencoder (RQ1)}\label{sec:ae}

Let $x_i=(x_{i,1},\dots,x_{i,T_i})$ be the standardised PER sequence of player $i$ and $m_{i,t}\in\{0,1\}$ the observed-indicator. Each season is embedded as a token $\bm{u}_{i,t}=\bm{W}_{\mathrm{in}}[x_{i,t}m_{i,t},\,m_{i,t}]^{\top}+\bm{p}_t$, with $\bm{p}_t$ a learned positional embedding. A learned classification token $\bm{u}_{i,0}$ is prepended and the sequence is processed by a Transformer encoder \citep{vaswani2017} with a key-padding mask that removes unobserved and padded positions from attention. The latent code is a linear projection of the final state of the classification token,
\begin{equation}
\bm{z}_i=\bm{W}_z\,\mathrm{Enc}(\bm{u}_{i,0},\dots,\bm{u}_{i,T})_0\in\mathbb{R}^{d},
\label{eq:latent}
\end{equation}
so that the latent dimensionality $d$ is decoupled from the model width. The decoder is a Transformer decoder whose queries are learned positional embeddings $\bm{q}_t$ attending through cross-attention to the single memory vector $\bm{W}_m\bm{z}_i$; a linear read-out gives $\hat{x}_{i,t}$. Training minimises the masked mean-squared error
\begin{equation}
\mathcal{L}=\frac{\sum_{i}\sum_{t=1}^{T} m_{i,t}\,(x_{i,t}-\hat{x}_{i,t})^{2}}{\sum_{i}\sum_{t=1}^{T} m_{i,t}},
\label{eq:loss}
\end{equation}
in which gaps and padding contribute nothing; no cluster information enters training. Table~\ref{tab:aeconfig} lists every architectural and optimisation setting. Latent dimensionalities $d\in\{8,16,32,64,128\}$ were trained with three seeds each ($d=64$ with five); the seed governs initialisation and mini-batch order, and deterministic algorithms were enforced. Held-out RMSE in PER units is compared with three baselines on the same observations: the development-set global mean, each player's own career mean and a player-specific least-squares linear trend. The prespecified criterion for retaining the autoencoder is a held-out RMSE below that of the linear trend.

\begin{table}[t]\centering\small
\caption{Transformer autoencoder: architecture, initialisation, optimisation and validation settings.}
\label{tab:aeconfig}
\begin{tabular}{ll}\toprule
Component & Setting \\ \midrule
Input per season & $[x_{i,t}m_{i,t},\,m_{i,t}]$ (standardised PER, observed indicator) \\
Model width / heads / feed-forward & 64 / 4 / 128 \\
Encoder / decoder layers & 2 / 2 (post-norm, ReLU, dropout 0.1) \\
Positional encoding & learned, season index (encoder) and output position (decoder) \\
Latent dimensionality $d$ & 8, 16, 32, 64, 128 (linear projection of the CLS state) \\
Initialisation & Xavier-uniform weights, zero biases, CLS $\sim\mathcal N(0,0.02^2)$ \\
Loss & masked MSE, Eq.~\eqref{eq:loss} \\
Optimiser & AdamW, learning rate $10^{-3}$, weight decay $10^{-4}$ \\
Batch size / gradient clipping & 64 / global norm 1.0 \\
Schedule & halve learning rate after 15 epochs without validation improvement \\
Epochs / early stopping & at most 400; patience 40 on held-out masked MSE; best epoch restored \\
Data split & 80/20 by player, stratified by career-length tertile, split seed 2026 \\
Seeds & 0--4 ($d=64$), 0--2 (other $d$); deterministic algorithms enforced \\
Parameters & 179,073 ($d=64$) \\
Hardware / time & 2 CPU cores; 8 min per model ($d=64$) \\
\bottomrule\end{tabular}
\end{table}

\subsection{Single-run UMAP+HDBSCAN and its reproducibility (RQ2)}\label{sec:singlerun}

The design of the earlier version, standardisation of the latent codes, UMAP \citep{mcinnes2018} and HDBSCAN \citep{campello2013,mcinnes2017hdbscan}, was re-implemented with a full grid (UMAP dimension $\{2,3,5,10\}$, neighbours $\{10,15,20,30,50\}$, minimum distance $\{0,0.1\}$; HDBSCAN minimum cluster size $\{15,20,25,30,40,50\}$, minimum samples $\{5,10,15\}$; \nGrid{} configurations, Euclidean distance, excess-of-mass extraction), evaluated on the development set only. Three selection rules were compared: (a) the prespecified rule, maximum DBCV \citep{moulavi2014} subject to at most 25\% noise and at least three clusters; (b) maximum Silhouette under the same constraints; (c) the composite score of the earlier version (Silhouette 30\%, noise share 30\%, capped Calinski--Harabasz 20\%, a step function of the number of clusters 20\%). Selected configurations were applied to the held-out players with UMAP's transform and HDBSCAN's approximate prediction, and refitted on all players. Reproducibility was measured directly: for four representative configurations the pipeline was rerun with ten UMAP random states for each of three autoencoder seeds, with excess-of-mass and leaf extraction, and every pair of runs was compared by the adjusted Rand index \citep[ARI;][]{hubert1985}, recording also the number of clusters, the noise share and the size of the largest cluster.

\subsection{Consensus clustering (RQ3)}\label{sec:consensus}

Because single runs proved irreproducible, the reported partitions are resampling-based consensus partitions \citep{monti2003,fred2005}. An ensemble member is an autoencoder seed (five for $d=64$), a random 80\% subsample of players and a random state; for each member and each $k\in\{2,\dots,12\}$ the subsample is partitioned by $k$-means (five initialisations) in the standardised latent space or, in the UMAP variant, in a five-dimensional UMAP embedding refitted on the subsample. With \nMembersAe{} members for the five-seed autoencoder ensemble and \nMembersSumm{} for single-embedding representations, the co-association matrix $C_k$ records for every pair of players the share of members containing both in which they were assigned together; the consensus partition for a given $k$ is the average-linkage agglomeration of $1-C_k$ into $k$ groups. Each $k$ is described by PAC, the share of pairs with $0.1<C_k<0.9$ \citep{senbabaoglu2014}, and by the consensus Silhouette on the distance $1-C_k$. The rule for $k$ selects the $k\ge3$ that maximises the consensus Silhouette. We had initially specified the PAC minimum; PAC decreased almost monotonically in $k$ for every representation (Appendix~\ref{app:k}) and could not select $k$, so the rule was changed before any cluster profile was examined, and both curves are reported. A player's stability is the mean co-association with the other members of its consensus cluster; players below 0.5 are reported as unstable. The same protocol with HDBSCAN in place of $k$-means was also run.

\subsection{Baseline representations and common criteria (RQ4)}\label{sec:baselines}

The protocol was applied unchanged to (i) nine summary statistics (career mean, peak, standard deviation, valid seasons, span, linear slope, first and last season PER, relative position of the peak, share of seasons at or above 15), standardised; (ii) the trajectory resampled by linear interpolation onto ten equally spaced points of normalised career time plus log career length, standardised (``resampled''); (iii) its first six principal components (\pcaVar\% of variance); (iv) $k$-means with dynamic time warping \citep{sakoe1978,petitjean2011,tavenard2020} on the variable-length standardised sequences (``DTW''); and (v) Ward linkage on the summary statistics. Every partition is scored on the same criteria: reproducibility (PAC, consensus Silhouette, share of stable players); trajectory-shape homogeneity, the share of variance of the resampled-trajectory matrix explained by cluster membership, $R^2_{\mathrm{shape}}=1-\mathrm{SS}_{\mathrm{within}}/\mathrm{SS}_{\mathrm{total}}$, a space no method optimised directly; and external validity, Cram\'er's $V$ and adjusted mutual information \citep{vinh2010} with position, draft category, All-Star and All-NBA selection, and the Kruskal--Wallis $\varepsilon^2$ \citep{kruskal1952,tomczak2014} for career WS, mean BPM and career games. Internal indices in each method's own space are reported in Appendix~\ref{app:internal} but not used across representations.

\subsection{Tests for discrete structure and a continuous description}\label{sec:structure_methods}

Reproducibility is necessary but not sufficient for the existence of types, because $k$-means partitions even a unimodal cloud stably \citep{senbabaoglu2014}. Four analyses address this. (i) The gap statistic \citep{tibshirani2001} compares $\log W_k$, the within-cluster dispersion of $k$-means, with its expectation under reference data without structure, for $k=1,\dots,12$ and $B=20$ draws from a uniform distribution over the principal-component bounding box and from a multivariate normal with the sample covariance; $k$ is chosen by the rule of \citet{tibshirani2001}. (ii) Hartigan's dip test of unimodality \citep{hartigan1985} is applied to the first three principal components of each representation and to individual summary statistics. (iii) The consensus protocol itself is applied to three draws from a multivariate normal with the mean and covariance of the real features and to three draws from a Gaussian copula that preserves every real marginal and the rank correlations but imposes a unimodal dependence structure. (iv) To separate level from shape, the protocol and the dip test are applied to a shape-only representation (each resampled trajectory minus its own mean, length excluded), the representation closest to the functional-data analyses of \citet{wakim2014}. As a constructive alternative to types, principal components of the resampled trajectory in PER units are related to level, slope, peak timing and length, and a linear discriminant on these coordinates is used to measure how much of the typology they reproduce.

\subsection{Sensitivity, labelling and software (RQ5)}\label{sec:sens_methods}

The full pipeline was repeated for latent dimensionalities 8, 16, 32 and 128, for a single autoencoder seed, unstandardised codes and a UMAP stage; and both the autoencoder pipeline and the summary-statistic typology were repeated for the six alternative sample specifications of Table~\ref{tab:flow}, for gap deletion and for BPM in place of PER. Agreement with the primary partition is measured by ARI on common players. Cluster labels follow a rule fixed before the results were examined: a level term from the cluster median of career-mean PER (below 12, 12--18, above 18), a duration term from the median of valid seasons (at most 4, 5--9, at least 10) and a shape qualifier only when the median slope exceeds 0.25 PER per season in absolute value. Named players are the three nearest to each cluster medoid, plus well-known players whose membership is discussed only as illustration. Analyses used Python 3.11 with PyTorch \citep{paszke2019}, scikit-learn \citep{pedregosa2011}, umap-learn, hdbscan \citep{mcinnes2017hdbscan} and tslearn \citep{tavenard2020} on two CPU cores; all seeds are fixed in the released code, which also regenerates every number in this paper from the result files.
